t
\documentclass[12pt]{article}
\usepackage[margin=1in]{geometry}
\usepackage{hyperref}
\usepackage{enumitem}
\usepackage{amsmath, amssymb}
\usepackage{graphicx}
\usepackage{parskip}
\usepackage[T1]{fontenc}
\usepackage[utf8]{inputenc}
\setlength{\parskip}{0.8em}

\title{\textbf{Projects Guide}\\[0.5em]
Natural Language Processing in Python --- HSE University}
\author{Assoc.\ Prof.\ Soroosh Shalileh}
\date{Spring Semester (Modules 3--4)}

\begin{document}
\maketitle

\section*{Overview}
This guide outlines expectations and evaluation criteria for the two major course projects in the \emph{Natural Language Processing in Python} course. Both projects are designed to foster research thinking, technical depth, and clear scientific communication.

\begin{itemize}[leftmargin=*]
  \item \textbf{Mini-Project (Module 3)}: Student-designed midterm project focused on foundational NLP methods.
  \item \textbf{Final Project (Module 4)}: Student-designed research or application project using advanced neural/transformer-based methods.
\end{itemize}

\section{Mini-Project (Module 3)}
\subsection*{Objective}
The Mini-Project consolidates fundamental NLP concepts covered in the first module of the course. It should combine preprocessing, representation, and modeling steps to form an end-to-end pipeline.

\subsection*{Structure}
\begin{itemize}[leftmargin=*]
  \item Work individually or in pairs.
  \item Select one of the guided tracks below or propose a similar variant.
  \item Duration: 2-4 weeks (Weeks 6--8).
  \item Deliverables: GitHub repository, short paper (2--3 pages), 5-minute presentation.
\end{itemize}

\subsection*{Suggested Tracks (Idea)}
\begin{enumerate}[leftmargin=*]
  \item \textbf{Classical NLP Pipeline:} Tokenization $\rightarrow$ TF--IDF $\rightarrow$ Logistic Regression or Feedforward Neural Network.\\
        \emph{Goal:} Compare preprocessing choices and analyze their impact on performance.\\
        \emph{Dataset:} IMDb reviews, AG News, or a Russian sentiment corpus.

  \item \textbf{Embedding-Based Classifier:} Use word2vec or GloVe embeddings and train Feedforward Neural Network classifier.\\
        \emph{Goal:} Compare pretrained vs.\ locally trained embeddings on small data.\\
        \emph{Dataset:} News or Wikipedia paragraphs.

  \item \textbf{Sequence Modeling (optional advanced):} Train an more complex (not being taught during the lecture)  on character- or word-level language modeling.\\
        \emph{Goal:} Explore overfitting and perplexity as a diagnostic tool.

  \item \textbf{Quantitative Study:} Analyze tokenizer or language model errors across languages or domains.\\
        \emph{Goal:} Demonstrate ability to collect and interpret linguistic data.
\end{enumerate}

\subsection*{Grading (100 points)}
\begin{tabular}{ll}
\textbf{Component} & \textbf{Points} \\
\hline
Technical correctness \& implementation & 30 \\
Experimental design \& analysis & 25 \\
Insight and interpretation & 20 \\
Code organization \& reproducibility & 15 \\
Presentation \& clarity of report & 10 \\
\end{tabular}

\subsection*{Timeline}
\begin{tabular}{lll}
\textbf{Week} & \textbf{Milestone} & \textbf{Deliverable}\\
\hline
6 & Topic selection & 1-paragraph plan approved by instructor \\
7 & Preliminary results & short check-in during lab \\
8 & Submission \& presentation & repository + short paper \\
\end{tabular}

\section{Final Project (Module 4)}
\subsection*{Objective}
The Final Project is a capstone research or applied project, designed and executed by students. It should demonstrate mastery of modern NLP methods, including transformers, retrieval, or alignment approaches (if time permits).

\subsection*{Structure}
\begin{itemize}[leftmargin=*]
  \item Work in teams of 1--3 students (2 recommended).
  \item Duration: Weeks 3--6 (6 weeks).
  \item Proposal approval required (Week 8 deadline). Sent an email to sshalileh@hse.ru, indicate title, dataset(s) and objective in one paragraph. 
  \item Deliverables: GitHub repository, 4--6 page paper, 10-minute demo presentation.
\end{itemize}

\subsection*{Example Project Ideas}
\begin{enumerate}[leftmargin=*]
  \item \textbf{Domain-Specific Fine-Tuning:} Adapt BERT or RoBERTa to a domain corpus (finance, legal, health). Evaluate with domain-specific metrics.
  \item \textbf{Cross-Lingual Transfer:} Compare multilingual transformers (mBERT, XLM-R) for NER or sentiment across languages.
  \item \textbf{Retrieval-Augmented QA:} Implement a small RAG pipeline using dense or sparse retrievers; evaluate latency vs.\ accuracy.
  \item \textbf{Instruction-Tuned Models:} Perform small-scale SFT or DPO on a compact model; design custom evaluation rubric for helpfulness.
  \item \textbf{Interpretability Study:} Analyze token-level saliency or probing on pretrained LMs; visualize hidden representations.
  \item \textbf{Bias and Fairness Analysis:} Examine demographic or topical bias in model outputs using template-based generation.
  \item \textbf{Low-Resource Adaptation:} Explore parameter-efficient fine-tuning (LoRA, prefix-tuning) for small corpora.
  \item \textbf{NLP Systems Project:} Build an applied pipeline (chatbot, text summarizer, domain QA) integrating modern APIs or models.
\end{enumerate}


\subsection*{Grading Rubric (100 points)}
\begin{tabular}{ll}
\textbf{Component} & \textbf{Points} \\
\hline
Problem definition \& clarity & 10 \\
Data \& experimental setup & 20 \\
Technical correctness \& methodology & 25 \\
Analysis \& interpretation & 25 \\
Writing \& presentation quality & 15 \\
Originality \& insight & 5 \\
\end{tabular}

% \subsection*{Timeline}
% \begin{tabular}{lll}
% \textbf{Week} & \textbf{Milestone} & \textbf{Deliverable}\\
% \hline
% 11 & Brainstorming session & group discussions during lab \\
% 12 & Proposal submission & 1-page PDF proposal \\
% 14 & Progress update & short in-class demo \\
% 16 & Final presentation & slides + 6--8 page paper + repo \\
% \end{tabular}

\section*{General Notes}
\begin{itemize}[leftmargin=*]
  \item Reports should follow the ACL two-column format (LaTeX or Overleaf template).
  \item Code must be reproducible and accompanied by a clear \texttt{README}.
  \item Datasets must be publicly available or properly licensed.
  \item Collaborations across teams are permitted for discussion but not code sharing.
  \item Outstanding projects may be invited for research paper submissions or lab collaboration.
\end{itemize}

\end{document}
